\section{Convexity of Logistic Regression [20 points]}
In this problem, we will study a function that arises in logistic regression. 
No prior knowledge of logistic regression is required. Everything needed to solve the problem is defined below, and the main goal is to reason about convexity.

Consider a binary classification problem with labels $y_i \in \{0,1\}$ and inputs $x_i \in \mathbb{R}^d$. For this problem, define the logistic regression model directly by
\begin{equation*}
    P(y_i=1 \mid x_i,w)=\sigma(w^\top x_i), \qquad \sigma(z)=\frac{1}{1+\exp(-z)}.
\end{equation*} 

\subsection{Convex Optimization}
The corresponding log-likelihood is 
\begin{equation*}
    \mathcal{L}(w) = \log P(y|\mathbf{X}, w) = \sum_{i=1}^n[y_i w^\top x_i - \log(1+ \exp(w^\top x_i))].
\end{equation*}
Our goal is to find the weight vector $w$ that maximizes this likelihood. In this question, you will prove that $\mathcal{L}$ is indeed a concave function (and hence, the negative conditional log likelihood is a convex function). 

\begin{enumerate}
    \item \textbf{[3 points]} A real-valued function $f:S \rightarrow \mathcal{R}$ defined on a convex set S, is said to be \textit{convex} if $$f(tx_{1}+(1-t)x_{2})\leq tf(x_{1})+(1-t)f(x_{2}), \forall x_{1},x_{2}\in S,\forall t\in [0,1].$$
    Show that a linear combination of $n$ convex functions, $f_1, f_2,...,f_n$, $\sum_{i=1}^n a_if_i(x)$ is also a convex function  $\forall a_i \in R^+$. 


\begin{tcolorbox}[fit,height=7cm,width=.95\textwidth,blank,borderline={1pt}{-2pt},nobeforeafter]
%solution

\end{tcolorbox}

\clearpage
\item \textbf{[1 point]} Show that a linear combination of $n$ concave functions, $f_1, f_2,...,f_n$, $\sum_{i=1}^n a_if_i(x)$ is also a concave function  $\forall a_i \in R^+$. Recall that if a function $f(x)$ is convex, then $-f(x)$ is concave. (You can use the result from part (1).) 
    

\begin{tcolorbox}[fit,height=7cm,width=.95\textwidth,blank,borderline={1pt}{-2pt},nobeforeafter]
%solution

\end{tcolorbox}
  
\item \textbf{[3 points]} Another property of twice differentiable convex functions is that the second derivative is non-negative. Using this property, show that $f (x) = \log (1 + \exp x)$ is a convex function. Note that this property is both sufficient and necessary once we assume that $f$ is twice differentiable: i.e., (if $f''(x)$ exists, then $f''(x)\ge 0 \iff  f$ is convex).
    
\begin{tcolorbox}[fit,height=7cm,width=.95\textwidth,blank,borderline={1pt}{-2pt},nobeforeafter]

\end{tcolorbox}
\clearpage
\item \textbf{[3 points]} Let $f_i:\mathcal{S} \rightarrow \mathcal{R}$ for $i = 1, \ldots, n$ be a set of convex functions.
Is $f(x) = \max_i f_i(x)$ also convex?
If yes, prove it. 
If not, provide a counterexample.  


\begin{tcolorbox}[fit,height=6cm,width=.95\textwidth,blank,borderline={1pt}{-2pt},nobeforeafter]
%solution

\end{tcolorbox}

\item \textbf{[8 points]} Show that the log likelihood of \textit{Logistic Regression} is a concave function. You may use the fact that if $f$ and $g$ are both convex and twice differentiable, and if $g$ is non-decreasing, then $g \circ f$ is convex.
    
    
\begin{tcolorbox}[fit,height=12cm,width=.95\textwidth,blank,borderline={1pt}{-2pt},nobeforeafter]
%solution

\end{tcolorbox}

\newpage

\item \textbf{[2 points]} Unfortunately, for this model, we cannot derive a closed-form solution with MLE. However, now that we have proven that $\mathcal{L}$ is concave, are there other methods we can use to solve for $w$? Is any method guaranteed to find the ``best" solution, where ``best" means achieving the lowest loss? In a few short words, explain why or why not.

\begin{tcolorbox}[fit,height=6cm,width=.95\textwidth,blank,borderline={1pt}{-2pt},nobeforeafter]
%solution

\end{tcolorbox}
  
\end{enumerate}
